\scanpage{92}

\[(90)\qquad X_i^*=\deleted{\unsure{X_j}}X_i\setminus\bigcup_{j<i}X_j\]
Les $X_i^*$ \uline{non vides} forment alors une partition de $X$. Les $X_i^*$ sont les ``strates ouvertes''.

On introduit parfois des axiomes supplémentaires :

(Strat 4) (fidélité) $\forall i\in I$, $X_i^*\neq\varnothing$.

Cet axiome a l'inconvénient de ne pas être stable par \deleted{\unsure{la}} restriction à un ouvert de $X$. Il est parfois complété par l'axiome plus fort, encore moins stable

(Strat 4 bis) (Canonicité de $I$) les $X_i^*$ sont $0$-connexes, i.e. connexes non vides,

(Strat 5) \inserted{(densité) $\forall i\in I$,} $X_i^*$ dense dans $X_i$, \inserted{i.e. $X_i=\overline{X_i^*}$;} i.e. $j<i\Rightarrow X_j$ \uline{rare} dans $X_i$.
\margin{moyennant (Strat 5), la condition de fidélité équivaut à $X_i\neq\varnothing$ $\forall i\in I$ (car $X_i=\overline{X_i^*}$ donc $X_i\neq\varnothing\Rightarrow X_i^*\neq\varnothing$)}

Cet axiome a l'avantage d'être stable par restriction à un ouvert. Mais on ne \unsure{supposera} (Strat 4) (voire Strat 4 bis) ou (Strat 5) que par mention expresse.

\deleted{\unsure{Stratifications de fidélité implique|Cet axiome de fidélité implique}}
On a pour tt $i\in I$
\[
(91)\qquad X_i=\bigcup_{\substack{j\in\bar i\\\text{i.e. }j\leq i}}X_j^*
\]
donc on aura, si $i,j\in J$
\margin{91}
